\documentclass[11pt]{article}

\usepackage[a4paper,margin=1in]{geometry}
\usepackage{amsmath,amssymb,amsthm,mathtools}
\usepackage{microtype}
\usepackage[hidelinks]{hyperref}

\newcommand{\HH}{\mathcal H}
\newcommand{\E}{\mathbb E}
\newcommand{\Prob}{\mathbb P}
\newcommand{\inner}[2]{\left\langle #1,#2\right\rangle_{\HH}}
\newcommand{\norm}[1]{\left\lVert #1\right\rVert_{\HH}}
\newcommand{\pos}[1]{\left[#1\right]_+}

% Proof headings: use these instead of \paragraph inside proofs.
\newcommand{\proofstep}[1]{%
  \par\smallskip
  \noindent\textbf{#1}\enspace\ignorespaces
}

\newtheorem{assumption}{Assumption}
\newtheorem{lemma}{Lemma}
\newtheorem{theorem}{Theorem}
\newtheorem{remark}{Remark}

\title{Energy-Normalized RKHS Pearson Classification:\\
Threshold and Separation Guarantees}
\author{}
\date{}

\begin{document}
\maketitle

\section{Pearson discrepancy and energy normalization}

Let $\HH$ be a real separable RKHS with measurable feature
map $\phi$. For a reference distribution $P_i$, define
\[
\mu_i=\E_{P_i}\phi(X),
\qquad
T_i=\E_{P_i}[\phi(X)\otimes\phi(X)],
\qquad
A_i=T_i+\lambda I,
\]
where $\lambda>0$ and the outer-product operator is defined by
\[
(u\otimes v)f=\inner{v}{f}u.
\]
For an incoming distribution $Q$, write
\[
\mu_Q=\E_Q\phi(Z),
\qquad
\Delta_i=\mu_Q-\mu_i.
\]

We use the regularized RKHS Pearson-type discrepancy
\begin{equation}
D_i(Q)=\frac12\inner{\Delta_i}{A_i^{-1}\Delta_i}.
\label{eq:pearson}
\end{equation}
Its connection to Pearson is the quadratic variational form
\begin{equation}
D_i(Q)
=
\sup_{f\in\HH}
\left\{
\E_Q f-\E_{P_i}f
-\frac12\left(\E_{P_i}f^2+\lambda\norm{f}^2\right)
\right\}.
\label{eq:variational}
\end{equation}
Indeed, the expression being maximized equals
\[
\inner{f}{\Delta_i}-\frac12\inner{f}{A_i f}
=
D_i(Q)
-\frac12
\norm{A_i^{1/2}(f-A_i^{-1}\Delta_i)}^2.
\]

Let
\[
v_i=A_i^{-1/2}\mu_Q,
\qquad
w_i=A_i^{-1/2}\mu_i.
\]
Then
\[
D_i(Q)=\frac12\norm{v_i-w_i}^2.
\]
Define the total energy and normalized score by
\begin{align}
E_i(Q)
&=
\inner{\mu_Q}{A_i^{-1}\mu_Q}
+\inner{\mu_i}{A_i^{-1}\mu_i},
\label{eq:energy}
\\
S_i(Q)
&=
\frac{D_i(Q)}{E_i(Q)}
=
\frac{\norm{v_i-w_i}^2}
{2(\norm{v_i}^2+\norm{w_i}^2)}.
\label{eq:score}
\end{align}
Thus, normalization preserves the reference second-moment
geometry and measures discrepancy relative to total
embedding energy. It defines a normalized Pearson-type score,
not the classical Pearson divergence.

If $E_i(Q)=0$, define $S_i(Q)=0$.
Since
\[
\norm{v_i-w_i}^2
\le2(\norm{v_i}^2+\norm{w_i}^2),
\]
we have $0\le S_i(Q)\le1$.

\section{Assumptions and empirical quantities}

\begin{assumption}[Kernel and sampling]
\label{ass:sampling}
The positive-semidefinite kernel is fixed before sampling
and satisfies
\[
k(x,y)\ge0,
\qquad
k(x,x)\le K
\]
for all $x,y$, where $K>0$.
The two windows are independent and satisfy
\[
X_1,\ldots,X_n\overset{\mathrm{iid}}{\sim}P_i,
\qquad
Z_1,\ldots,Z_n\overset{\mathrm{iid}}{\sim}Q.
\]
\end{assumption}

These kernel assumptions imply
\begin{equation}
\norm{\phi(x)}\le\sqrt K,
\qquad
\norm{\phi(x)-\phi(y)}\le\sqrt{2K}.
\label{eq:feature-bounds}
\end{equation}
Indeed,
\[
\norm{\phi(x)-\phi(y)}^2
=k(x,x)+k(y,y)-2k(x,y)\le2K.
\]
The second bound uses pointwise nonnegativity of the kernel.

Define
\[
\widehat\mu_i=\frac1n\sum_{j=1}^n\phi(X_j),
\qquad
\widehat\mu_Q=\frac1n\sum_{j=1}^n\phi(Z_j),
\]
and
\[
\widehat T_i
=\frac1n\sum_{j=1}^n\phi(X_j)\otimes\phi(X_j),
\qquad
\widehat A_i=\widehat T_i+\lambda I.
\]
The empirical quantities are
\begin{align}
\widehat D_i
&=
\frac12
\inner{\widehat\mu_Q-\widehat\mu_i}
{\widehat A_i^{-1}(\widehat\mu_Q-\widehat\mu_i)},
\\
\widehat E_i
&=
\inner{\widehat\mu_Q}
{\widehat A_i^{-1}\widehat\mu_Q}
+
\inner{\widehat\mu_i}
{\widehat A_i^{-1}\widehat\mu_i},
\\
\widehat S_i
&=\frac{\widehat D_i}{\widehat E_i}.
\end{align}
Set $\widehat S_i=0$ if $\widehat E_i=0$.
The same squared-norm inequality as above gives
$0\le\widehat S_i\le1$.

The incoming window is accepted as compatible with reference
$i$ when $\widehat S_i\le\eta$.

\begin{assumption}[Reference embedding lower bound]
\label{ass:energy}
A known constant $m_0>0$ satisfies
\[
\norm{\mu_i}\ge m_0.
\]
\end{assumption}

This assumption controls the normalization denominator
under the null. It is separate from the Pearson separation
assumption introduced below.

\section{Deterministic energy bounds}

\subsection{Basic operator inequalities}

For every $f\in\HH$, Cauchy--Schwarz gives
\[
\begin{aligned}
\inner{f}{\widehat T_i f}
&=\frac1n\sum_{j=1}^n
|\inner{f}{\phi(X_j)}|^2\\
&\le\frac1n\sum_{j=1}^n
\norm{f}^2\norm{\phi(X_j)}^2\\
&\le K\norm{f}^2.
\end{aligned}
\]
Hence
\[
0\preceq\widehat T_i\preceq KI,
\qquad
\lambda I\preceq\widehat A_i\preceq(K+\lambda)I.
\]
Inversion reverses the operator order:
\begin{equation}
\frac1{K+\lambda}I
\preceq\widehat A_i^{-1}
\preceq\frac1\lambda I.
\label{eq:inverse-bounds}
\end{equation}
The same inequalities hold for $T_i$ and $A_i$.

To justify the inverse-order rule, suppose $A\succeq C$,
where both operators are bounded, self-adjoint, and bounded
below by a positive multiple of $I$. Then
\[
C^{-1/2}AC^{-1/2}\succeq I,
\]
and consequently
\[
A^{-1}
=
C^{-1/2}(C^{-1/2}AC^{-1/2})^{-1}C^{-1/2}
\preceq C^{-1}.
\]
This argument remains valid in an infinite-dimensional RKHS.

\subsection{Energy lower bound}

\begin{lemma}[Midpoint decomposition]
\label{lem:midpoint}
Let
\[
\widehat c_i=\frac{\widehat\mu_Q+\widehat\mu_i}{2}.
\]
Then
\begin{align}
\widehat E_i
&=
\widehat D_i
+2\inner{\widehat c_i}
{\widehat A_i^{-1}\widehat c_i},
\label{eq:midpoint}\\
\widehat E_i
&\ge
\widehat D_i+\frac{2\norm{\widehat c_i}^2}{K+\lambda}.
\label{eq:energy-lower}
\end{align}
\end{lemma}

\begin{proof}
Write $x=\widehat\mu_Q$, $y=\widehat\mu_i$, and
$B=\widehat A_i^{-1}$.
Since $B$ is self-adjoint,
\[
\widehat D_i
=
\frac12\inner{x}{Bx}
+\frac12\inner{y}{By}
-\inner{x}{By}.
\]
On the other hand,
\[
\begin{aligned}
2\inner{\widehat c_i}{B\widehat c_i}
&=\frac12\inner{x+y}{B(x+y)}\\
&=\frac12\inner{x}{Bx}
+\frac12\inner{y}{By}
+\inner{x}{By}.
\end{aligned}
\]
Adding cancels the cross terms and proves
\eqref{eq:midpoint}.

Using \eqref{eq:inverse-bounds},
\[
\inner{\widehat c_i}{B\widehat c_i}
\ge\frac{\norm{\widehat c_i}^2}{K+\lambda},
\]
which proves \eqref{eq:energy-lower}.
\end{proof}

\subsection{Energy upper bound}

\begin{lemma}[Uniform energy upper bound]
\label{lem:upper}
For every incoming distribution $Q$,
\[
E_i(Q)\le M_\lambda,
\qquad
\widehat E_i\le M_\lambda,
\]
where
\begin{equation}
M_\lambda=\frac K\lambda+\frac K{K+\lambda}.
\label{eq:M}
\end{equation}
\end{lemma}

\begin{proof}
\proofstep{Step 1: Bound the incoming energy.}
Every empirical mean embedding satisfies
\[
\norm{\frac1n\sum_j\phi(X_j)}
\le\frac1n\sum_j\norm{\phi(X_j)}
\le\sqrt K.
\]
Consequently,
\[
\inner{\widehat\mu_Q}
{\widehat A_i^{-1}\widehat\mu_Q}
\le\frac{\norm{\widehat\mu_Q}^2}{\lambda}
\le\frac K\lambda.
\]

\proofstep{Step 2: Use the reference covariance identity.}
Write $y=\widehat\mu_i$.
Expanding the empirical covariance gives
\[
\widehat T_i-y\otimes y
=
\frac1n\sum_j
(\phi(X_j)-y)\otimes(\phi(X_j)-y)
\succeq0.
\]
Therefore,
\[
\widehat A_i\succeq\lambda I+y\otimes y,
\qquad
\widehat A_i^{-1}
\preceq(\lambda I+y\otimes y)^{-1}.
\]

\proofstep{Step 3: Evaluate the rank-one inverse along $y$.}
By definition,
\[
(y\otimes y)y=\norm{y}^2y.
\]
It follows that
\[
(\lambda I+y\otimes y)
\left(\frac{y}{\lambda+\norm{y}^2}\right)=y.
\]
Hence
\[
(\lambda I+y\otimes y)^{-1}y
=\frac{y}{\lambda+\norm{y}^2},
\]
including when $y=0$.
Thus,
\[
\begin{aligned}
\inner{y}{\widehat A_i^{-1}y}
&\le
\inner{y}{(\lambda I+y\otimes y)^{-1}y}\\
&=\frac{\norm{y}^2}{\lambda+\norm{y}^2}\\
&\le\frac K{\lambda+K}.
\end{aligned}
\]
The last inequality follows because
$t\mapsto t/(\lambda+t)$ is increasing for $t\ge0$.

Notice that $y$ is an eigenvector of $y\otimes y$ when
$y\ne0$; we do not assume it is an eigenvector of
$\widehat T_i$.

\proofstep{Step 4: Add the two terms.}
We obtain
\[
\widehat E_i
\le\frac K\lambda+\frac K{\lambda+K}.
\]
For the population result, repeat the argument using
\[
\norm{\mu_Q},\norm{\mu_i}\le\sqrt K
\]
and
\[
T_i-\mu_i\otimes\mu_i
=
\E_{P_i}\big[
(\phi(X)-\mu_i)\otimes(\phi(X)-\mu_i)
\big]\succeq0.
\]
\end{proof}

\section{Null guarantee}

Assume $Q=P_i$. For $\rho_0\in(0,1)$, define
\begin{align}
U_0
&=\frac K{n\lambda}
\left(1+\sqrt{\log\frac2{\rho_0}}\right)^2,
\\
a_0
&=\sqrt{\frac K{2n}}
\left(1+\sqrt{\log\frac2{\rho_0}}\right),
\\
H_0
&=\frac{2[m_0-a_0]_+^2}{K+\lambda},
\\
\eta_0
&=\frac{U_0}{U_0+H_0}.
\label{eq:eta0}
\end{align}

\begin{theorem}[Acceptance under the null]
\label{thm:null}
Under Assumptions~\ref{ass:sampling} and~\ref{ass:energy},
\[
Q=P_i
\quad\Longrightarrow\quad
\Prob\{\widehat S_i\le\eta_0\}\ge1-\rho_0.
\]
Moreover, $\eta_0<1$ whenever $m_0>a_0$.
\end{theorem}

\begin{proof}
\proofstep{Step 1: Expected mean-difference error.}
Write $\mu=\mu_i$ and
\[
F=\norm{\widehat\mu_Q-\widehat\mu_i}.
\]
Under the null, independence and centering give
\[
\begin{aligned}
\E F^2
&=
\E\norm{(\widehat\mu_Q-\mu)-(\widehat\mu_i-\mu)}^2\\
&=
\E\norm{\widehat\mu_Q-\mu}^2
+\E\norm{\widehat\mu_i-\mu}^2\\
&=\frac2n\left(\E k(X,X)-\norm{\mu}^2\right)
\le\frac{2K}{n}.
\end{aligned}
\]
Therefore,
\[
\E F\le\sqrt{\frac{2K}{n}}.
\]

\proofstep{Step 2: Concentration of the mean difference.}
Replacing one of the $2n$ observations changes $F$
by at most $\sqrt{2K}/n$. The sum of squared
bounded-difference constants is
\[
2n\left(\frac{\sqrt{2K}}n\right)^2=\frac{4K}{n}.
\]
One-sided McDiarmid concentration yields
\[
\Prob\{F-\E F\ge t\}
\le\exp\left(-\frac{nt^2}{2K}\right).
\]
Taking
\[
t=\sqrt{\frac{2K}{n}\log\frac2{\rho_0}}
\]
gives, with probability at least $1-\rho_0/2$,
\[
F\le\sqrt{\frac{2K}{n}}
\left(1+\sqrt{\log\frac2{\rho_0}}\right).
\]
Hence
\[
\widehat D_i\le\frac{F^2}{2\lambda}\le U_0.
\]

\proofstep{Step 3: Concentration of the midpoint.}
Under the null, $\widehat c_i$ is the empirical mean
of $2n$ iid observations from $P_i$.
For
\[
G=\norm{\widehat c_i-\mu},
\]
we have
\[
\E G\le\sqrt{\E G^2}\le\sqrt{\frac K{2n}}.
\]
Replacing one observation changes $G$ by at most
$\sqrt{2K}/(2n)$, so the sum of squared constants is
\[
2n\left(\frac{\sqrt{2K}}{2n}\right)^2=\frac Kn.
\]
Therefore,
\[
\Prob\{G-\E G\ge t\}
\le\exp\left(-\frac{2nt^2}{K}\right).
\]
Taking
\[
t=\sqrt{\frac K{2n}\log\frac2{\rho_0}}
\]
gives $G\le a_0$ with probability at least $1-\rho_0/2$.

\proofstep{Step 4: Combine numerator and denominator bounds.}
By a union bound, both concentration events hold with
probability at least $1-\rho_0$.
Their independence is not required.

On their intersection,
\[
\norm{\widehat c_i}
\ge[\norm{\mu}-a_0]_+
\ge[m_0-a_0]_+.
\]
Lemma~\ref{lem:midpoint} gives
\[
\widehat E_i\ge\widehat D_i+H_0.
\]
If $H_0>0$, then
\[
\widehat S_i
\le\frac{\widehat D_i}{\widehat D_i+H_0}
\le\frac{U_0}{U_0+H_0},
\]
because $d\mapsto d/(d+H_0)$ is increasing.
If $H_0=0$, then $\eta_0=1$ and the result follows
from $\widehat S_i\le1$.
\end{proof}

\section{Standard separation guarantee}

The standard proof controls the full mean-embedding errors
and the operator error. It is conservative because the
operator norm controls every feature-space direction,
including directions unnecessary for a particular pair.

\begin{assumption}[Pearson separation]
\label{ass:separation}
For the fixed pair $(Q,P_i)$,
\[
D_i(Q)\ge\delta>0.
\]
\end{assumption}

At the population level, Lemma~\ref{lem:upper} immediately
gives
\begin{equation}
\boxed{S_i(Q)\ge\frac{\delta}{M_\lambda}.}
\label{eq:population-separation}
\end{equation}
Indeed, $D_i(Q)>0$ implies $E_i(Q)>0$, and
$E_i(Q)\le M_\lambda$.

\subsection{Joint concentration}

For $\rho_1\in(0,1)$, define
\begin{align}
b_n(\rho_1)
&=\sqrt{\frac Kn}
\left(1+\sqrt{\log\frac3{\rho_1}}\right),
\label{eq:b}
\\
e_n(\rho_1)
&=\frac K{\sqrt n}
\left(1+\sqrt{2\log\frac3{\rho_1}}\right).
\label{eq:e}
\end{align}

\begin{lemma}[Mean and operator concentration]
\label{lem:concentration}
With probability at least $1-\rho_1$, all three
inequalities hold:
\[
\norm{\widehat\mu_Q-\mu_Q}\le b_n(\rho_1),
\qquad
\norm{\widehat\mu_i-\mu_i}\le b_n(\rho_1),
\]
and
\[
\|\widehat T_i-T_i\|_{\mathrm{op}}\le e_n(\rho_1).
\]
\end{lemma}

\begin{proof}
\proofstep{Step 1: Concentration of one mean embedding.}
For an iid empirical mean $\widehat\mu$ with population
mean $\mu$,
\[
\E\norm{\widehat\mu-\mu}^2
=\frac1n\left(\E k(X,X)-\norm{\mu}^2\right)
\le\frac Kn.
\]
Thus
\[
\E\norm{\widehat\mu-\mu}\le\sqrt{\frac Kn}.
\]
Replacing one observation changes the norm error by
at most $\sqrt{2K}/n$. McDiarmid gives
\[
\Prob\left\{
\norm{\widehat\mu-\mu}>\sqrt{\frac Kn}+t
\right\}
\le\exp\left(-\frac{nt^2}{K}\right).
\]
Take
\[
t=\sqrt{\frac Kn\log\frac3{\rho_1}}.
\]
This proves each mean bound with failure probability
at most $\rho_1/3$.

\proofstep{Step 2: Expected operator error.}
Let $V=\phi(X)\otimes\phi(X)$.
In the Hilbert space of Hilbert--Schmidt operators,
\[
\|V\|_{\mathrm{HS}}=\norm{\phi(X)}^2\le K.
\]
Independence and centering imply
\[
\E\|\widehat T_i-T_i\|_{\mathrm{HS}}^2
=\frac1n\E\|V-\E V\|_{\mathrm{HS}}^2
\le\frac{K^2}{n}.
\]
Consequently,
\[
\E\|\widehat T_i-T_i\|_{\mathrm{HS}}
\le\frac K{\sqrt n}.
\]

\proofstep{Step 3: Concentration of the operator error.}
Replacing one observation changes the Hilbert--Schmidt
norm error by at most
\[
\frac1n
\|\phi(x)\otimes\phi(x)-\phi(x')\otimes\phi(x')\|_{\mathrm{HS}}
\le\frac{2K}{n}.
\]
Thus,
\[
\Prob\left\{
\|\widehat T_i-T_i\|_{\mathrm{HS}}
>\frac K{\sqrt n}+t
\right\}
\le\exp\left(-\frac{nt^2}{2K^2}\right).
\]
Take
\[
t=K\sqrt{\frac2n\log\frac3{\rho_1}}
\]
and use
\[
\|\widehat T_i-T_i\|_{\mathrm{op}}
\le\|\widehat T_i-T_i\|_{\mathrm{HS}}.
\]
This proves the operator bound with failure probability
at most $\rho_1/3$.

\proofstep{Step 4: Combine the events.}
A union bound over the two mean bounds and the operator
bound gives total failure probability at most $\rho_1$.
\end{proof}

\subsection{Empirical score lower bound}

Define
\[
s
=\sqrt{\frac2\lambda}\,b_n(\rho_1),
\qquad
r=\frac{e_n(\rho_1)}{\lambda},
\]
and
\begin{equation}
L(\delta,\rho_1)
=
\frac{[\sqrt\delta-s]_+^2}
{M_\lambda(1+r)}.
\label{eq:Lstd}
\end{equation}

\begin{theorem}[Alternative guarantee]
\label{thm:standard}
Under Assumptions~\ref{ass:sampling}
and~\ref{ass:separation},
\[
\Prob\{\widehat S_i\ge
L(\delta,\rho_1)\}\ge1-\rho_1.
\]
Consequently, if $\eta<L(\delta,\rho_1)$,
then
\[
\Prob\{\widehat S_i>\eta\}\ge1-\rho_1.
\]
\end{theorem}

\begin{proof}
Work on the event of Lemma~\ref{lem:concentration}.

\proofstep{Step 1: Bound the mean-difference error.}
Set
\[
\widehat\Delta_i=\widehat\mu_Q-\widehat\mu_i.
\]
The triangle inequality gives
\[
\begin{aligned}
\norm{\widehat\Delta_i-\Delta_i}
&\le
\norm{\widehat\mu_Q-\mu_Q}
+\norm{\widehat\mu_i-\mu_i}\\
&\le2b_n(\rho_1).
\end{aligned}
\]

\proofstep{Step 2: Compare empirical and population metrics.}
The operator error bound implies
\[
\widehat A_i\preceq A_i+e_n(\rho_1)I.
\]
Since $A_i\succeq\lambda I$,
\[
e_n(\rho_1)I
\preceq\frac{e_n(\rho_1)}{\lambda}A_i.
\]
Hence
\[
\widehat A_i
\preceq(1+r)A_i.
\]
Inverting gives
\[
\widehat A_i^{-1}
\succeq\frac{A_i^{-1}}{1+r}.
\]

\proofstep{Step 3: Apply the reverse triangle inequality.}
Define
\[
\|f\|_{A_i^{-1}}
=\sqrt{\inner{f}{A_i^{-1}f}}.
\]
Because $A_i^{-1}\preceq\lambda^{-1}I$,
\[
\|\widehat\Delta_i-\Delta_i\|_{A_i^{-1}}
\le\frac{2b_n(\rho_1)}{\sqrt\lambda}.
\]
Therefore,
\[
\begin{aligned}
\|\widehat\Delta_i\|_{A_i^{-1}}
&\ge
\left[
\|\Delta_i\|_{A_i^{-1}}
-\|\widehat\Delta_i-\Delta_i\|_{A_i^{-1}}
\right]_+\\
&\ge
\left[
\sqrt{2D_i(Q)}
-\frac{2b_n(\rho_1)}{\sqrt\lambda}
\right]_+.
\end{aligned}
\]

\proofstep{Step 4: Lower-bound empirical Pearson.}
Combining the previous inequalities,
\[
\begin{aligned}
\widehat D_i
&\ge
\frac{\|\widehat\Delta_i\|_{A_i^{-1}}^2}
{2(1+r)}\\
&\ge
\frac{[\sqrt{D_i(Q)}-s]_+^2}
{1+r}\\
&\ge
\frac{[\sqrt\delta-s]_+^2}
{1+r}.
\end{aligned}
\]

\proofstep{Step 5: Normalize using the energy upper bound.}
If this lower bound is positive, then $\widehat E_i>0$.
Lemma~\ref{lem:upper} gives
\[
\widehat S_i
=\frac{\widehat D_i}{\widehat E_i}
\ge\frac{\widehat D_i}{M_\lambda}
\ge L(\delta,\rho_1).
\]
If the bound is zero, the claim follows from
$\widehat S_i\ge0$, including the convention for
$\widehat E_i=0$.
\end{proof}

\section{Nonempty threshold interval}

Combining the null and standard alternative guarantees,
a certified threshold satisfies
\[
\boxed{
\eta_0\le\eta<L(\delta,\rho_1).
}
\]
This interval is nonempty if and only if
\[
\eta_0<
\frac{[\sqrt\delta-s]_+^2}
{M_\lambda(1+r)}.
\]
Since $\eta_0>0$, the positive part must be strictly positive.
Multiplying by the denominator and taking square roots gives
\[
\sqrt\delta-s
>
\sqrt{M_\lambda(1+r)\eta_0}.
\]
Equivalently,
\begin{equation}
\boxed{
\delta>
\delta_{\mathrm{crit}}
:=
\left[
s
+\sqrt{M_\lambda(1+r)\eta_0}
\right]^2.
}
\label{eq:critical-standard}
\end{equation}
This is a necessary and sufficient condition for the
displayed certified interval to be nonempty, not a necessary
condition for successful classification.

For $\rho_0=\rho_1=\rho$, write
\[
U=\frac K{n\lambda}
\left(1+\sqrt{\log\frac2\rho}\right)^2,
\]
and
\[
H=\frac2{K+\lambda}
\left[
m_0-\sqrt{\frac K{2n}}
\left(1+\sqrt{\log\frac2\rho}\right)
\right]_+^2.
\]
Then
\[
\eta_0=\frac{U}{U+H},
\]
and the explicit condition is
\[
\boxed{
\begin{aligned}
\delta>
\Bigg[
&\sqrt{\frac{2K}{n\lambda}}
\left(1+\sqrt{\log\frac3\rho}\right)
\\
&+
\sqrt{
M_\lambda
\left[
1+\frac K{\lambda\sqrt n}
\left(1+\sqrt{2\log\frac3\rho}\right)
\right]
\frac{U}{U+H}
}
\Bigg]^2.
\end{aligned}
}
\]
The null proof combines two concentration events, explaining
$\log(2/\rho)$. The standard alternative combines three,
explaining $\log(3/\rho)$.

If $m_0\le a_0$, then $\eta_0=1$ and no useful interval
can be certified. For fixed positive $m_0,K,\lambda,\rho$,
as $n\to\infty$,
\[
U=O(n^{-1}),
\qquad
H\longrightarrow\frac{2m_0^2}{K+\lambda}>0.
\]
Consequently,
\[
\eta_0=O(n^{-1}),
\qquad
\delta_{\mathrm{crit}}=O(n^{-1}).
\]

\appendix

\section{Alternative directional proof}

The standard proof controls errors in all feature-space
directions. To certify separation of a fixed pair, one can
instead control a single informative population direction.
This gives smaller error terms without changing the score
or its implementation.

Define
\[
s_{\mathrm{dir}}
=\sqrt{\frac K{n\lambda}\log\frac2{\rho_1}},
\qquad
r_{\mathrm{dir}}
=\frac K\lambda
\sqrt{\frac1{2n}\log\frac2{\rho_1}},
\]
and
\begin{equation}
L_{\mathrm{dir}}(\delta,\rho_1)
=
\frac{[\sqrt\delta-s_{\mathrm{dir}}]_+^2}
{M_\lambda(1+r_{\mathrm{dir}})}.
\label{eq:Ldir}
\end{equation}

\begin{theorem}[Directional alternative guarantee]
\label{thm:directional}
Under Assumptions~\ref{ass:sampling}
and~\ref{ass:separation},
\[
\Prob\{\widehat S_i\ge
L_{\mathrm{dir}}(\delta,\rho_1)\}\ge1-\rho_1.
\]
Consequently, if $\eta<L_{\mathrm{dir}}(\delta,\rho_1)$,
then
\[
\Prob\{\widehat S_i>\eta\}\ge1-\rho_1.
\]
\end{theorem}

\begin{proof}
Write $D=D_i(Q)$, $A=A_i$, and $\Delta=\Delta_i$.

\proofstep{Step 1: Choose a fixed population direction.}
Since $D>0$, define
\[
u=\frac{A^{-1}\Delta}{\sqrt{2D}}.
\]
This direction depends on the population distributions,
not on the sampled windows. It is only a proof device;
the classifier does not need to compute it.

Using $2D=\inner{\Delta}{A^{-1}\Delta}$ gives
\[
\inner{u}{\Delta}=\sqrt{2D},
\qquad
\inner{u}{Au}=1.
\]
Since $A\succeq\lambda I$,
\[
1=\inner{u}{Au}\ge\lambda\norm{u}^2,
\qquad
\norm{u}^2\le\frac1\lambda.
\]

\proofstep{Step 2: Restrict the empirical variational problem.}
Define
\[
\widehat d=\inner{u}{\widehat\mu_Q-\widehat\mu_i},
\qquad
\widehat q=\inner{u}{\widehat A_i u}.
\]
Restricting the empirical version of
\eqref{eq:variational} to $f=tu$, $t\ge0$, gives
\[
\widehat D_i
\ge
\sup_{t\ge0}
\left(t\widehat d-\frac{t^2}{2}\widehat q\right).
\]
Since $\widehat q>0$, the maximizing value is
$t=[\widehat d]_+/\widehat q$. Hence
\begin{equation}
\widehat D_i\ge\frac{[\widehat d]_+^2}{2\widehat q}.
\label{eq:directional-variational}
\end{equation}

\proofstep{Step 3: Concentrate the directional mean.}
Let
\[
h(x)=\inner{u}{\phi(x)}.
\]
Using \eqref{eq:feature-bounds},
\[
|h(x)-h(y)|
\le\norm{u}\norm{\phi(x)-\phi(y)}
\le\sqrt{\frac{2K}{\lambda}}.
\]
Thus the range of $h$ has length at most
$\sqrt{2K/\lambda}$.

Moreover,
\[
\widehat d
=\frac1n\sum_jh(Z_j)-\frac1n\sum_jh(X_j),
\qquad
\E\widehat d=\sqrt{2D}.
\]
The sum of squared range lengths of the $2n$ independent
summands is at most
\[
2n\left(\frac1n\sqrt{\frac{2K}{\lambda}}\right)^2
=\frac{4K}{n\lambda}.
\]
One-sided Hoeffding concentration gives
\[
\Prob\{\widehat d<\sqrt{2D}-t\}
\le\exp\left(-\frac{n\lambda t^2}{2K}\right).
\]
Taking $t=\sqrt2\,s_{\mathrm{dir}}$ yields
\[
\widehat d\ge\sqrt{2D}-\sqrt2\,s_{\mathrm{dir}}
\]
with probability at least $1-\rho_1/2$.

\proofstep{Step 4: Concentrate the directional quadratic term.}
We have
\[
\widehat q
=\frac1n\sum_jh(X_j)^2+\lambda\norm{u}^2,
\qquad
\E\widehat q=\inner{u}{Au}=1.
\]
Also,
\[
0\le h(X_j)^2
\le\norm{u}^2\norm{\phi(X_j)}^2
\le\frac K\lambda.
\]
Hoeffding therefore gives
\[
\Prob\{\widehat q>1+t\}
\le\exp\left(-\frac{2n\lambda^2t^2}{K^2}\right).
\]
Taking $t=r_{\mathrm{dir}}$ yields
\[
\widehat q\le1+r_{\mathrm{dir}}
\]
with probability at least $1-\rho_1/2$.

\proofstep{Step 5: Combine the two scalar bounds.}
By a union bound, both events hold with probability
at least $1-\rho_1$. Their independence is not required.
On their intersection, \eqref{eq:directional-variational}
gives
\[
\begin{aligned}
\widehat D_i
&\ge
\frac{[\sqrt{2D}-\sqrt2\,s_{\mathrm{dir}}]_+^2}
{2(1+r_{\mathrm{dir}})}\\
&=
\frac{[\sqrt D-s_{\mathrm{dir}}]_+^2}
{1+r_{\mathrm{dir}}}\\
&\ge
\frac{[\sqrt\delta-s_{\mathrm{dir}}]_+^2}
{1+r_{\mathrm{dir}}}.
\end{aligned}
\]

\proofstep{Step 6: Normalize.}
When this bound is positive, $\widehat E_i>0$ and
Lemma~\ref{lem:upper} gives
\[
\widehat S_i
\ge\frac{\widehat D_i}{M_\lambda}
\ge L_{\mathrm{dir}}(\delta,\rho_1).
\]
When the bound is zero, the claim follows from
$\widehat S_i\ge0$.
\end{proof}

\subsection{Threshold condition and comparison}

The null guarantee is unchanged.
The directional certified interval is
\[
\eta_0\le\eta<L_{\mathrm{dir}}(\delta,\rho_1).
\]
As in the standard analysis, it is nonempty precisely when
\begin{equation}
\boxed{
\delta>
\delta_{\mathrm{crit,dir}}
:=
\left[
s_{\mathrm{dir}}
+\sqrt{M_\lambda(1+r_{\mathrm{dir}})\eta_0}
\right]^2.
}
\label{eq:critical-directional}
\end{equation}

For the same parameters,
\[
s_{\mathrm{dir}}
=
\sqrt{\frac K{n\lambda}}
\sqrt{\log\frac2{\rho_1}},
\]
whereas
\[
s
=
\sqrt{\frac{2K}{n\lambda}}
\left(1+\sqrt{\log\frac3{\rho_1}}\right).
\]
Likewise,
\[
r_{\mathrm{dir}}
=
\frac K{\lambda\sqrt n}
\sqrt{\frac12\log\frac2{\rho_1}},
\]
whereas
\[
r
=
\frac K{\lambda\sqrt n}
\left(1+\sqrt{2\log\frac3{\rho_1}}\right).
\]
Therefore,
\[
s_{\mathrm{dir}}<s,
\qquad
r_{\mathrm{dir}}<r.
\]
It follows that
\[
L_{\mathrm{dir}}(\delta,\rho_1)
\ge L(\delta,\rho_1),
\qquad
\delta_{\mathrm{crit,dir}}
<
\delta_{\mathrm{crit}}.
\]
The lower-bound inequality is strict whenever
$L_{\mathrm{dir}}(\delta,\rho_1)>0$.

\end{document}